\documentclass[10pt, letterpaper]{article}

\usepackage{../../../../template/template}

% ------------------------
% Impostazioni documento
% ------------------------

\setDocTitle{Verbale esterno 07/04/2026}
\setDocVersion{-}
\setDocData{07/04/2026}
\setDocIndex{ve\_2026\_04\_07}
\setDocState{2}
\setDocRecipient{2}

\begin{document}

\makeTitlePage

\newpage

\tableofcontents

\newpage

% -------------------------------------------------
% Informazioni riunione
% -------------------------------------------------
\makeMeetingInfoTable{07/04/2026}{16:00}{16:30}{via \textit{Microsoft Teams$_G$}}

% ------------------------
% Partecipanti
% ------------------------
\addParticipant{Filippo}{Venzo}{Programmatore}{P}
\addParticipant{Valeria}{Baleanu}{Programmatore}{P}
\addParticipant{Luca}{Granziero}{Amministratore}{P}
\addParticipant{Christian}{Libralato}{Responsabile}{P}
\addParticipant{Francesco}{Pasqual}{Programmatore}{P}
\addParticipant{Leonardo}{Pellizzon}{Programmatore}{P}
\addParticipant{Giuseppe}{De Fina}{Verificatore}{P}
\makeMeetingParticipantsTable

% -------------------------------------------------
% Ordine del giorno
% -------------------------------------------------
\section{Ordine del giorno}
\begin{itemize}
      \item Stato di Avanzamento del Progetto;
      \item preparazione e Struttura della Presentazione;
      \item modellazione Architetturale C4;
      \item pianificazione delle Prossime Scadenze e Collaudi.
\end{itemize}

% -------------------------------------------------
% Approfondimento
% -------------------------------------------------
\section{Approfondimento}

\subsection*{Stato di Avanzamento del Progetto}
È stato presentato un aggiornamento sullo stato del progetto, evidenziando lo stato ancora incompleto della parte documentale e di alcune funzionalità come la gestione delle notifiche.
La maggior parte del sistema è tuttavia già implementata e funzionante; il gruppo ha confermato che sta
lavorando per completare quanto più possibile in vista delle prossime scadenze.

In particolare è emerso che la logica di invio delle notifiche tramite websocket è già implementata sia lato \textit{front-end$_G$} che \textit{back-end$_G$}, ma manca la visualizzazione grafica
che dovrebbe essere completata in serata.
Inoltre il gruppo ha confermato di aver configurato il droplet, acquistato il dominio viewforlife.me tramite GitHub Developer Pack e installato il certificato SSL con certbot, mettendo online il sistema.

\subsection*{Preparazione e Struttura della Presentazione}
Successivamente è stata discussa la struttura della presentazione, definendo le funzionalità principali da mostrare, la struttura delle slide e le modalità di esposizione, con particolare attenzione
agli aspetti tecnici e ai flussi di utilizzo. La Proponente ha consigliato di strutturare la presentazione con un approccio top-down, partendo dall'architettura, mostrando i flussi di configurazione,
gestione eventi/allarmi e analytics, associando a ciascun flusso uno schema e uno screenshot, il tutto evitando i dettagli troppo tecnici.

\subsection*{Modellazione Architetturale C4}
Sono stati presentati i primi tre livelli del modello C4, discutendo quali servizi e componenti includere. È stata notata la mancanza del componente esterno \textit{Groq$_G$}, utilizzato per la
generazione dei suggerimenti, ed è stato chiarito che è sufficiente rappresentare i primi due livelli nella presentazione.

\subsection*{Pianificazione delle Prossime Scadenze e Collaudi}
Sono infine state confermate le tempistiche per il pre-collaudo, la presentazione al marketing e il collaudo finale, con possibilità di variare, 
per quanto possibile, le date in base alle esigenze del gruppo.
% -------------------------------------------------
% Decisioni
% -------------------------------------------------

% -------------------------------------------------
% Azioni
% -------------------------------------------------
\addTodo{\noOne}{Inviare la bozza della presentazione a Sciacco entro le 12:00 del 08/04 per una revisione e correzione.}{C. Libralato}{08/04/2026}
\addTodo{\noOne}{Anticipare il link del sito viewforlife.me e fornire due credenziali a Sciacco per permettergli di effettuare un accesso e una verifica}{F. Venzo}{08/04/2026}
\addTodo{\noOne}{Aggiornare il modello C4 inserendo il servizio \textit{Groq} come box separato}{F. Pasqual}{08/04/2026}
\addTodo{\noOne}{Aggiungere nella presentazione schemi a blocchi che rappresentino i tre flussi principali: configurazione iniziale, gestione eventi/allarmi, e generazione analytics con suggerimenti.}{V. Baleanu}{08/04/2026}

\makeTodoTable

\end{document}